\documentclass[11pt,a4paper]{article}
\usepackage[utf8]{inputenc}
\usepackage{graphicx}
\usepackage[left=2.5cm,top=3cm,right=2.5cm,bottom=3cm,bindingoffset=0.5cm]{geometry}
\usepackage{AEDMacros}
\usepackage{caratula}
\newlength{\miSangria}
\setlength{\miSangria}{2em}

% Asi pueden escribir nuevos comandos. 
% Este por ejemplo asegura q los nombres 
% que figuren con una tipografia diferenciada  
\newcommand{\Tipo}[1]{\mathsf{#1}}
% la sintaxis es \newcommand{\nombreDeLaMacro}[cantidadDeParametros]{Lo que va ser remplazado por el macro} 
\newcommand{\norm}[1]{\vert#1\vert}


\begin{document}

\section{Procedimiento Canjear Millas}
\subsection{Renombres de Tipos}

\noindent topeDePromocion \textbf{ES}: $\mathbb{R}$
\vspace{0.5cm}

\noindent enum Categoría \textbf{ES} \{Bronce, Plata,Oro, Platino  \}

\vspace{0.5cm}

\noindent Transacciones \textbf{ES} Struct$\struct{
		UsuarioOrigen: \Tipo{\Z}, \\
		UsuarioDestino: \Tipo{\Z},
		TipoDeTransaccion: \Tipo{\Z},
		Millas: \Tipo{\R}, \\
		NombrePromo: \Tipo{\str},
		FactorPromo: \Tipo{\R}
	}$

\noindent \comentario{Tipo de transacciones es algo que yo pongo a mano creo, ejemplo 0:Acumulacion, 1: Transferencia, 2: Canjeo}

\subsection{TAD y Observadores}

\begin{tad}{Airmiles}



	\obs{Usuarios}{\dict{id:\Tipo{\Z}}{CategoriaActual : Categoria}}

	\noindent \obs{HistorialDeTransacciones}{\seq{Transacciones}}
	\\
	\texttt{obs} Promociones : dict \{ \\
	\hspace*{1.5em} NombrePromocion: $\str$, \\
	\hspace*{1.5em} $\struct{FactorDePromocion:\Tipo{\R},TopeDePromocion:\Tipo{\R}, MillasRestantes:\Tipo{\R}}$ \\
	\}
	\vspace{0.5cm}

	\begin{proc}{CanjearMillas}
		{\Inout AM: AirMiles, \In MillasCanjeadas : \Tipo{\R}, \In IdUsuario : \Tipo{\Z}}
		{}
		\requiereLargo{ AM=AM_{0} \yLuego MillasCanjeadas > 0 \land IdUsuario \in claves(AM.Usuarios)}
		\aseguraLargo{
			\IfThenElse{ MillasDisponibles(AM_{0}.Historial, IdUsuario) \geq millasCanjeadas}
			{ \\TransaccionCanje(AM_{0}.Historial, IdUsuario, MillasCanjeadas)}
			{ \\AM.Historial =  AM_{0}.Historial} \\
			\yLuego AM.Usuarios = AM_{0}.Usuarios \\
			\yLuego AM.Promociones = AM_{0}.Promociones
		}
	\end{proc}
	\\ \\
	\comentario{Predicado para transaccion de canje. \\ La idea es tener el agregado de la transaccion 'filtrado' porque en la tupla de transaccion \\hay campos que dependen del tipo de transaccion y creo que da claridad.}
	\\
	\predLargo{TransaccionCanje}{historialTransacciones: \seq{Transacciones},idUsuario : \Tipo{\Z},\\ millasCanjeadas : \Tipo{\R}}
	{\tupla{idUsuario,-1,2,millasCanjeadas,'NoPromo',1.0} \in historialTransacciones}

	\vspace{0.5cm}
	\comentario{Auxiliar para obtener las millas disponibles de un usuario. \\ Recorre el historial de transacciones y de acuerdo al TipoDeTransaccion, al IdUsuario \\ y las Millas del registro, suma o resta con un acumulador}
	\\

	\auxLargo{MillasDisponibles}
	{IdUsuario : \Tipo{\Z}, Historial : \seq{Transacciones}}
	{\Tipo{\Z}}
	{res = \sum_{i=0}^{\lvert Historial \rvert - 1} \IfThenElse{EsTransaccionDelUsuario(Historial[i], IdUsuario)}{\\\IfThenElse{EsTransaccionSumadora(Historial[i], IdUsuario)\\}{Historial[i].Millas}{-Historial[i].Millas}}{0}}

	\vspace{1.0cm}
	\predLargo{EsTransaccionSumadora}{Trans: Transacciones,idUsuario : \Tipo{\Z}}
	{ (Trans.TipoDeTransaccion = 0) \lor \\(IdUsuario = Trans.UsuarioDestino \land Trans.TipoDeTransaccion = 1)}

	\comentario{En la primer condicion no valido el IdUsuario porque se supone que si no es una transferencia nunca deberia venir en el campo UsuarioDestino, y ya valide que al menos esta en uno de los 2.}

	\vspace{1.0cm}
	\pred{EsTransaccionDelUsuario}{Trans: Transacciones,idUsuario : \Tipo{\Z}}{\\Trans.UsuarioOrigen = idUsuario \oLuego Trans.UsuarioDestino = idUsuario}

	\comentario{Valido que el IdUsuario este en alguno de los 2 campos. Si no esta en ninguno, entonces no debo sumar ni restar nada.}


\end{tad}
\end{document}